\documentclass{article}
\usepackage{amsmath, amssymb, amsthm}

\setlength{\parindent}{0pt}

\newtheorem{claim}{Claim}

\begin{document}

\begin{claim}
    There are $15120$ ways to arrange BANANARAMA.
\end{claim}

\begin{proof}
    The word BANANARAMA consists of $10$ letters, with $5$ A's, $2$ N's, and one each of B, R, and M.
    Since identical letters are indistinguishable, the number of distinct arrangements is
    \[
    \frac{10!}{2! \cdot 5!} = 15120
    \]

    \medskip

    Hence there are $15120$ ways to arrange BANANARAMA.
\end{proof}

\begin{claim}
    There are $\frac{72!}{12! \cdot 24! \cdot 36!}$ different paths from point $(0, 0, 0)$ to 
    $(12, 24, 36)$, assuming every step affects only one coordinate.
\end{claim}

\begin{proof}
    The entire path consists of $72$ steps, with $12$ steps in the $x$ direction, $24$ steps in the $y$ 
    direction, and $36$ steps in the $z$ direction. Since steps in the same coordinate are 
    indistinguishable, the number of distinct paths is
    \[
    \frac{72!}{12! \cdot 24! \cdot 36!}
    \]

    \medskip

    Hence there are $\frac{72!}{12! \cdot 24! \cdot 36!}$ different paths from point $(0, 0, 0)$ to 
    $(12, 24, 36)$, assuming every step affects only one coordinate.
\end{proof}

\begin{claim}
    There are $\frac{(2n)!}{2^n \cdot n!}$ different ways to pair up $2n$ students.
\end{claim}

\begin{proof}
    There are $(2n)!$ linear orderings of $2n$ students. Reading each ordering as consecutive pairs 
    produces a pairing. However, this procedure overcounts each pairing. Firstly, the two students 
    within each pair can be interchanged, giving a factor of $2^n$. Secondly, the $n$ pairs 
    themselves can be permuted arbitrarily, giving an additional factor of $n!$. Hence each pairing 
    is overcounted by $2^n \cdot n!$, so the number of distinct pairings is
    \[
    \frac{(2n)!}{2^n \cdot n!}
    \]

    \medskip

    Hence there are $\frac{(2n)!}{2^n \cdot n!}$ different ways to pair up $2n$ students.
\end{proof}

\begin{claim}
    There are $26{,}075{,}972{,}546$ solutions in $\mathbb{N}_0$ to the equation $x_1 + x_2 + x_3 + 
    \dots + x_8 = 100$.
\end{claim}

\begin{proof}
    The number of solutions in nonnegative integers corresponds to distributing $100$ indistinguishable 
    objects among $8$ distinguishable boxes. By the stars and bars theorem, the number of such 
    distributions is
    \[
    \binom{100 + 8 - 1}{8 - 1} = 26075972546
    \]

    \medskip

    Hence there are $26{,}075{,}972{,}546$ solutions in $\mathbb{N}_0$ to the equation $x_1 + x_2 + x_3 + 
    \dots + x_8 = 100$.
\end{proof}

\begin{claim}
    There are $2^n$ different ways to choose $n$ out of $2n$ objects, given that $n$ of the $2n$ objects 
    are identical and the other $n$ are unique.
\end{claim}

\begin{proof}
    Since the identical objects are indistinguishable, a selection is completely determined by which unique 
    objects are chosen. If $k$ unique objects are selected, then the remaining $n - k$ selected objects must 
    be identical objects. Thus, for each $k$, there are $\binom{n}{k}$ distinct selections. Hence the total 
    number of distinct selections is
    \[
    \sum_{k=0}^{n}\binom{n}{k} = 2^n
    \]

    \medskip

    Hence there are $2^n$ different ways to choose $n$ out of $2n$ objects, given that $n$ of the $2n$ 
    objects are identical and the other $n$ are unique.
\end{proof}

\begin{claim}
    There are $2^{\frac{n(n+1)}{2}}$ undirected graphs with vertices $v_1, v_2, \dots, v_n$ if self-loops 
    are permitted.
\end{claim}

\begin{proof}
    First we represent a graph by a binary string in which each bit corresponds to one possible edge. A bit 
    is $1$ if the edge is present and $0$ otherwise. This correspondence is one-to-one, since every graph 
    determines a unique binary string, and every binary string determines a unique graph.

    \medskip

    Order the bits according to the difference in the indices of their endpoints. First encode the $n$ 
    self-loops, then, for each $k \in \{1, 2, \dots, n - 1\}$, encode the $n - k$ edges joining $v_i$ and 
    $v_{i+k}$ for $1 \le i \le n - k$. Hence there are
    \[
    n + (n - 1) + (n - 2) + \dots + 1 = \frac{n(n+1)}{2}
    \]
    bits in the encoding.

    \medskip

    Since each bit can be $0$ or $1$ independently, there are
    \[
    2^{\frac{n(n+1)}{2}}
    \]
    distinct binary strings of this form. By the one-to-one correspondence above, this is exactly the number 
    of undirected graphs with vertices $v_1, v_2, \dots, v_n$ where self-loops are permitted.

    \medskip

    Hence there are $2^{\frac{n(n+1)}{2}}$ undirected graphs with vertices $v_1, v_2, \dots, v_n$ if 
    self-loops are permitted.
\end{proof}

\begin{claim}
    Assuming the ball only travels from left to right, there are $576$ distinct ways for it to bounce on 
    at most $3$ of $15$ sidewalk squares.
\end{claim}

\begin{proof}
    Since the ball only travels from left to right, a bounce pattern is uniquely determined by the set of 
    sidewalk squares in which the ball bounces on. Conversely, every choice of $k$ distinct sidewalk squares 
    determines exactly one such bounce pattern. This correspondence is one-to-one, hence the number of 
    distinct bounce patterns with exactly $k$ bounces is
    \[
    \binom{15}{k}
    \]
    Summing over all admissable values of $k$ gives
    \[
    \begin{aligned}
    \sum_{k=0}^{3}\binom{15}{k} &= \binom{15}{0} + \binom{15}{1} + \binom{15}{2} + \binom{15}{3} \\ 
    &= 576
    \end{aligned}
    \]

    \medskip

    Hence assuming the ball only travels from left to right, there are $576$ distinct ways for it to bounce 
    on at most $3$ of $15$ sidewalk squares.
\end{proof}

\begin{claim}
    There are $220$ work days avoided out of $300$ days if one skips days that are even-numbered, 
    multiples of $3$, and multiples of $5$.
\end{claim}

\begin{proof}
    Let $A$ be the set of days that are even-numbered, $B$ be the set of days that are multiples of $3$, 
    and $C$ be the set of days that are multiples of $5$. Then
    \[
    \begin{aligned}
    |A| &= \left\lfloor \frac{300}{2} \right\rfloor = 150 \\
    |B| &= \left\lfloor \frac{300}{3} \right\rfloor = 100 \\
    |C| &= \left\lfloor \frac{300}{5} \right\rfloor = 60
    \end{aligned}
    \]

    \medskip

    However, these three sets are not disjoint, since, for example, $6$ is in both $A$ and $B$. Hence the 
    total number of avoided work days is given by the inclusion-exclusion principle,
    \[
    |A \cup B \cup C| = |A| + |B| + |C| - |A \cap B| - |A \cap C| - |B \cap C| + |A \cap B \cap C|
    \]
    We compute the intersections as follows:
    \[
    \begin{aligned}
    |A \cap B| &= \left\lfloor \frac{300}{6} \right\rfloor = 50 \\
    |A \cap C| &= \left\lfloor \frac{300}{10} \right\rfloor = 30 \\
    |B \cap C| &= \left\lfloor \frac{300}{15} \right\rfloor = 20 \\
    |A \cap B \cap C| &= \left\lfloor \frac{300}{30} \right\rfloor = 10
    \end{aligned}
    \]
    Hence the total number of work days avoided is
    \[
    |A \cup B \cup C| = 150 + 100 + 60 - 50 - 30 - 20 + 10 = 220
    \]

    \medskip

    Hence there are $220$ avoided work days out of $300$ days if one skips days that are even-numbered, 
    multiples of $3$, and multiples of $5$.
\end{proof}

\end{document}
